\textbf{Scale and domain.} One chain length ($N=16$), a two-dimensional lattice, a binary alphabet, and \rhval{data/enumeration/hp16/n_designable_conformations/mean} possible targets, about 90 of them per test set. Real inverse folding has 20 amino acids, continuous backbones and approximate oracles; no claim here transfers to it. The oracle being an exact table lookup means ``budget'' is a count of objective evaluations, not a measured cost. The equal budget counts only per-design oracle calls: the AR designer's training set comes from scoring every sequence in the space (all 65,536, with full enumeration of conformations) at fraction 1.0, a cost amortised over all targets and not charged to $K$. The data-fraction sweep (scoring about a tenth of the space) still reaches \rhval{analysis/data-fraction-analysis/hp16/frac0p1_succ_k100_mean/mean:3} at $K=100$, but it was not run at an explicitly matched total oracle cost.

\textbf{Statistics.} Five seeds; seeds change the split and the initialisation together. Paired $t$-tests on five values have little power, none is corrected for multiplicity, and ``no detectable difference'' is not equivalence. The per-regime results (Table~\ref{tab:regime}) were not tested.

\textbf{Baselines and tuning.} The simulated annealing, random search and contact heuristic are our reimplementations, not published code. SA was tuned on a nine-configuration grid with one proposal type; the selected $T_0=0.5$ was the lowest value of that grid, so we added six lower-temperature configurations per objective on DEV; none beat the selection score of the original choice (best extended score \rhval{analysis_tune/sa-tuning-summary/hp16/bin_ext_best_score/mean:3} against \rhval{analysis_tune/sa-tuning-summary/hp16/bin_orig_best_score/mean:3} for the binary objective, \rhval{analysis_tune/sa-tuning-summary/hp16/log_ext_best_score/mean:3} against \rhval{analysis_tune/sa-tuning-summary/hp16/log_orig_best_score/mean:3} for the graded one), so the choice was kept, but the grid is still coarse and the heuristic-initialised variants were not tuned separately; other move sets (e.g.\ H-count-preserving swaps), other objectives or parallel tempering were not tried. The graded-objective and heuristic-initialised SA variants were added after seeing the main results, so they are exploratory, and the finding that they match or beat the AR model at $K=1000$ is the most important caveat on the registered comparison. The AR designer was tuned on four configurations; epochs were not retuned for the augmentation ablations, whose step-matched control (Table~\ref{tab:aug}) shows that training length matters; larger models, other conditioning features or sampling strategies (e.g.\ beam or oracle-guided decoding) were not explored. The DEV partition is a single split with about 90 targets.

\textbf{Design choices that matter.} The objective $f$, the choice to select the best of $K$ candidates by the oracle, and the success criterion (unique ground state) all shape the numbers; the contact heuristic and its failure analysis suggest that the target set (conformations that are designable at all) favours compact, contact-rich patterns. The failure analysis of registered SA is based on the energy-gap metric rather than on a separate intervention isolating degeneracy as the cause; the graded-objective runs are consistent with that explanation but do not prove it.

\textbf{What a full study would add.} More chain lengths (the enumeration grows exponentially, so $N\le 20$--$24$ is the practical limit), 3D lattices, more seeds and a split-by-seed decomposition, tuned baselines of equal effort, learned-model-guided search, and a comparison with a structure predictor as an approximate oracle.
